\section{Joint optical--decoder training}
\label{sec:s:joint}

\paragraph{Shared optics.}
The jointly trained counterpart uses the same phase parameterization, initial phase, frozen random surround, vectorial relay, detector footprints and reconstruction basis as the model-only designs. The measurement matrix is assembled differentiably from the excitation and fixed detection footprints. During joint training, each sampled 32-direction excitation is normalized by its own specimen-field mean. Evaluation instead uses the full-scan dose normalization shared by every encoder. This difference is retained explicitly in the training protocol; model-only correlation losses are invariant to a uniform scale.

\paragraph{Regularized solve and network.}
For $M=32\times144=4608$ measurements and $N=36{,}864$ coefficients, the decoder first computes
\begin{equation}
\mathbf z=\mathbf m+\bA^\top(\bA\bA^\top+\lambda\mathsf I)^{-1}
(\by-\bA\mathbf m).
\end{equation}
Here $\mathbf m$ is a constant volume filled with the training set's mean intensity. A learned log-parameter sets $\lambda>0$, initialized from mean training noise variance divided by training voxel variance. The solve is performed in measurement space. The Gram matrix is formed in float32 with TF32 disabled, and its Cholesky factorization and triangular solve use float64. The result is projected back in float32. One such solve is followed by one residual 3D U-Net.

The network has three encoder widths, 16, 32 and 64. Each block contains two $3\times3\times3$ convolutions with batch normalization and ReLU. Two $2\times2\times2$ max-pooling stages reduce the spatial grid. Transposed convolutions of kernel and stride two upsample; lateral skip features are concatenated before the corresponding convolution blocks. A final $1\times1\times1$ convolution predicts a residual that is subtracted from the network input. The output has no nonnegative or box activation.

\paragraph{Training data and noise.}
Training uses the same 459 neuronal crops from stacks 8 and 9 of the fourth animal as the denoiser study. Stack 10 contributes 240 held-out crops, of which every fifth crop gives 48 validation volumes. Each validation uses two fixed direction subsets and fixed noise, shared by all optical designs. Test animals never enter these stages. Training augments volumes with independent flips along each axis and lateral transposition.

Supervised training generates measurements from the coarse reconstruction-basis model $\bar{\by}=\bA\bx$. The noise approximation is
\begin{equation}
\widetilde y_i=\bar y_i+
\sqrt{\max(\bar y_i,0)/\kappa_k+10^{-12}}\,\xi_i,
\quad \xi_i\sim\mathcal N(0,1),
\end{equation}
with $\kappa_k=\kappa_0\,256/32$ and per-crop exposure calibrated using the untrained plate. Test measurements use the common fine-grid specimen model and exact Poisson sampling. This preserves the sub-voxel specimen residual during evaluation rather than adding it to supervised training targets.

\paragraph{Optimization schedule.}
All supervised stages use batch size 16 and reconstruction MSE. The initial decoder is warmed up for 5000 Adam steps on the frozen untrained plate at learning rate $10^{-3}$. Plate and decoder are then updated jointly for 1200 plate steps. The plate learning rate starts at 0.05 and follows a cosine decay to 0.005 over the last 30\,\% of steps. Decoder learning rate is $3\times10^{-4}$; each joint step is followed by four decoder-only steps on fresh batches using the same direction subset and a detached operator. The phase outside the trained central region remains fixed.

The initial 1200-step run did not meet a prespecified plateau criterion: validation MSE changed by 3.2\,\% over its last 300 plate updates, above the 2\,\% threshold. Before test evaluation, a fixed continuation of 600 updates was specified, with phase learning rate 0.005, unchanged decoder recipe, fresh optimizer states and restored model parameters. The continuation's last-300-step change was 0.025\,\%. The selected phase checkpoint, at update 1500, minimizes validation MSE across the continuation checkpoints. Phase RMS change from initialization is 1.44 radians. These checks limit selection and verify training behavior; they are not a claim of global convergence.

\paragraph{Numerical checks.}
At initialization, differentiable assembly reproduces stored operator rows to relative error $8.5\times10^{-8}$ after accounting for the subset/full-scan normalization scalar. An adjoint check and three central finite-difference gradient projections pass; directional derivatives agree to 0.03--0.07\,\%. A short descent check confirms that the phase changes and the loss decreases. Cross-server reconstruction checks compare recorded image-quality and centroid metrics rather than requiring bit-identical floating-point builds. These checks precede transfer comparisons.

\paragraph{Decoder comparison after optical training.}
The phase is fixed before evaluation with TV, pretrained PnP or neuron-trained PnP. Each operator uses its validation-selected scalar parameters, while denoisers remain unchanged. For the feed-forward comparison, every plate receives 3000 additional decoder-only steps at learning rate $3\times10^{-4}$, with validation every 500 steps. The untrained and model-only plates start this stage from the same warmed decoder; the end-to-end plate starts from its selected jointly trained checkpoint. Thus the architecture and additional fine-tuning budget match, but the preceding decoder training histories do not. The model with lowest validation MSE is used for each plate.

Feed-forward outputs are retained without final clipping for PSNR, foreground SSIM and the implemented centroid metric. Axial-profile W1 uses their nonnegative part, as specified in \cref{sec:s:evaluation}. This differs from the final clipping of TV/PnP outputs and is shared across all plates within the feed-forward comparison.

\paragraph{Comparison scope and timing.}
The primary end-to-end comparison uses one optical initialization and one supervised training run. Its plate and decoder were optimized for MSE; the primary reported axial metric is W1. On the RTX 5080, decoder warm-up, joint training and continuation took about 4500, 11{,}150 and 5574 seconds, respectively; per-plate fine-tuning took about 2700 seconds. The approximately 20-minute model-only timing was measured on an RTX 4500 Ada. We do not infer a controlled runtime speedup from these different machines.

An interpretation amendment was recorded after the initial TV/pretrained-PnP and own-decoder results, before the neuron-trained-PnP and off-design evaluations. Final reporting retains signed differences under every method; it does not classify the coupled plate as unable to transfer. New Fisher/coherence feed-forward evaluations use the same warm-up checkpoint and 3000-step recipe. Additional acquisition conditions are reported only where complete evaluations exist.
